In \textit{Shaw v.\ Reno}, 509 U.S.\ 630 (1993), the Supreme Court struck down
a North Carolina congressional district that was ``so bizarre on its face''
that it appeared to be a racial classification.
The district at issue (NC-12 in the 1992 plan) was famously described as
resembling a ``bug splattered on a windshield''---a district that connected
Black-majority communities along a thin corridor spanning the entire state.
The Court's legal test was implicitly visual: a district shape that ``rationally
cannot be understood as anything other than an effort'' at race-based sorting
triggers strict scrutiny~\citep{shaw1993}.

The convex hull ratio (CH) operationalises this visual-irregularity test
mathematically. A district like NC-12 (1992 version) has a convex hull that
is many times larger than the district itself, because the convex hull fills
the rectangular envelope of the district while the district occupies only a
thin corridor. CH $= A(\text{district}) / A(\text{convex hull})$; a narrow
corridor district has CH $\approx 0.10$--$0.20$.

\subsection{Why CH Detects What PP and Reock Miss}

Consider a ``cross-shaped'' district: four arms extending north, south, east, west
from a central hub, with a narrow corridor connecting each arm.
\begin{itemize}
  \item \textbf{PP} measures perimeter relative to area. A cross shape has
    substantial perimeter (each arm has four sides), but if the arms are
    thick enough, PP may register moderately (PP $\approx 0.10$--$0.25$).
  \item \textbf{Reock} measures area relative to the MBC. The MBC of a cross
    is determined by the tips of opposite arms; for a cross with arm length $L$
    and width $w$, MBC radius $\approx L\sqrt{2}/2$.
    Reock $\approx 4Lw / (L^2 \pi / 2) = 8w/(\pi L)$.
    For $L = 10$, $w = 2$: Reock $\approx 0.51$---a moderate value that
    would not flag the district as non-compact.
  \item \textbf{CH} measures area relative to the convex hull, which for a cross
    is approximately the square $L \times L$.
    CH $\approx 4Lw / L^2 = 4w/L$.
    For $L = 10$, $w = 2$: CH $\approx 0.80$---still passing, but for a more
    elongated cross ($L = 10$, $w = 0.5$), CH $= 0.20$, correctly flagging the shape.
\end{itemize}
For tentacle districts with very thin arms, CH is the most sensitive metric.

\subsection{Bisect's Structural Guarantee}

Recursive graph bisection produces convex-like sub-regions at each step because
METIS-based graph cuts tend to separate well-connected (spatially clustered)
regions. We demonstrate that all four bisect algorithms produce CH $> 0.85$
on NC 2020 by construction---a guarantee not available from hand-drawn maps.
